\scanpage{186}
l'assertion précédente équivaut à celle-ci :

\uline{pour tout diviseur} $D$ \uline{sur} $C$ \uline{tel que} $a\notin\operatorname{supp}D$, \uline{on a} $D\sim 0$ \uline{dans} $C-a$.

Pour le prouver, introduisons le cône projetant projectif $\widehat C$, et $E=\widehat C-a$, qui est un fibré vectoriel sur $X$ [savoir celui associé de façon covariante à $\underline O_X(1)$]. Alors $D$ provient d'un diviseur $\widehat D$ sur $\widehat C$ tel que $a\notin\operatorname{supp}\widehat D$, et par suite le \deleted{support de $\widehat D$ a} diviseur $\deleted{\unsure{E}}\unsure{D'|D}$ sur $E'$ induit par $\widehat D$ est \uline{à support propre} \uline{sur $X$}. Nous allons prouver le

\uline{Lemme} Soit $X$ préschéma régulier, $E$ un fibré vectoriel \inserted{de rang $1$} sur $X$ ($E=\underline W(\underline E)$, où $\underline E$ est un faisceau inv. sur $X$) $D$ un diviseur sur $E$, tel que suppt
% archive sheet number: 185
